\documentclass[10pt,a4paper]{amsart}
\usepackage[margin=19mm]{geometry}
\usepackage{amsmath,amssymb,mathtools}
\usepackage[scr=rsfs,cal=euler]{mathalfa}
\usepackage{xfrac}
\usepackage[unicode,hidelinks,bookmarks=false]{hyperref}
\newcommand{\ie}{i.e.\ }
\newcommand{\cf}{cf.\ }
\newcommand{\resp}{respectively\ }
\newcommand{\Subsection}{\subsection}
\providecommand{\nobreakdash}{\nobreak\hbox{-}}
\setlength{\emergencystretch}{2em}
\multlinegap0pt
\usepackage{enumerate}
\usepackage{amssymb}
\usepackage{float}
\usepackage{tikz}
\usepackage{xcolor}
\usepackage{algorithm}\usepackage{algpseudocode}
\newtheorem{lem}{Lemma}
\title{Proof of a Conjecture on Young Tableaux with Walls}
\author{Zhicong Lin$^{\color{blue} \dag}$, Feihu Liu$^{\color{blue} \ddag}$, Jiahang Liu$^{\color{blue} \S}$, Jing Liu$^{\color{blue}\P}$, and Guoce Xin$^{\color{blue} \maltese}$}
\date{Source: arXiv:2601.09551v3 (CC BY 4.0)}
\begin{document}
\maketitle

\noindent\emph{Source and licence.} Proof of a Conjecture on Young Tableaux with Walls. Version arXiv:2601.09551v3 (CC BY 4.0), distributed by arXiv under the Creative Commons Attribution 4.0 International licence, http://creativecommons.org/licenses/by/4.0/, as recorded in the arXiv licence metadata for this exact version and shown on the abstract page https://arxiv.org/abs/2601.09551v3. Full licence notice: \texttt{licenses/arxiv-2601.09551-CC-BY-4.0.txt}.\par\smallskip

\noindent\textit{Editorial scope.} Counted unit: the article's lem Theorem-Omega-k-1 (source0.tex lines 741--751). The article's complete original proof of it is delivered with it (source0.tex lines 752--852, one \texttt{proof} environment of 101 lines). No mathematics is written, repaired or re-derived: the statement and the proof are the article's own lines, in the article's own notation, and no retained line is substituted or re-flowed.  Retained uncounted as the source-local context the proof needs: lines 691--693.  Nothing was omitted from any retained span, so the package carries no omission record.  No retained line contains a citation, so the wrapper carries no bibliography and no bibliography entry.  No retained line contains a figure environment, an image-inclusion command or drawing code, so no image is delivered and the package declares no asset dependency.  Editorial formatting only, in addition: an AMS \texttt{amsart} preamble that restates the article's own declarations and macros used by the retained lines, namely the environments \texttt{lem} on separate counters, together with the standard packages those lines need.  The retained environments print in this document as Lemma 1 in order of appearance, whereas the article prints the same items with the numbering of its own full document.  The complete native source of the article as distributed in the arXiv e-print for this exact version is bundled unchanged as \texttt{sources/192688/source0.tex}.
\begin{equation}\label{rec:omega}
\omega_{n,m,k}=\omega_{n-1,m+1,k}-(m-k+2)\omega_{n-1,m+1,k-1}-\omega_{n-2,m+1,k}
\end{equation}

\begin{lem}\label{Theorem-Omega-k-1}
For any $1\leq s\leq n$ and $(n,k,s)\neq (1,0,1)$, the sequence $\omega_{n,k-1,k}$ satisfies the following recursion:
\begin{align*}
\omega_{n,k-1,k}=\sum_{p=1}^{\left\lfloor \frac{s+1}{2}\right\rfloor} \sum_{q=1}^{s+2-2p} \alpha_s(p,q)\cdot \omega_{n-s-1,k+s-p,k+1-q}
-\sum_{p=1}^{\left\lceil \frac{s+1}{2}\right\rceil} \sum_{q=1}^{s+3-2p} \alpha_{s+1}(p,q)\cdot \omega_{n-s,k+s-p,k+1-q},
\end{align*}
where $p, q\geq 1$ are integers, and
\begin{align}\label{Equation-Alpha}
\alpha_{s}(p,q)=\frac{(-1)^{q-p+1}\cdot (s-1+q-p)!}{(s-q-2p+2)!\cdot (q-1)!\cdot 2^{q-1}\cdot (p-1)!}.
\end{align}
\end{lem}

\begin{proof}
The proof is by induction on $s$.
If $s=1$, then we have
\begin{align*}
\omega_{n,k-1,k}&=\alpha_1(1,1)\cdot \omega_{n-2,k,k}-\alpha_2(1,1)\omega_{n-1,k,k}-\alpha_2(1,2)\omega_{n-1,k,k-1}
\\ &=-\omega_{n-2,k,k}+\omega_{n-1,k,k}-\omega_{n-1,k,k-1}.
\end{align*}
This is consistent with the recurrence for $\omega_{n,k-1,k}$ in~\eqref{rec:omega}, hence the lemma holds for $s=1$.

For $s\ge 2$, assume inductively that the formula holds for $s-1$. We now prove it for $s$.

By the recurrence for $\omega_{n,m,k}$ in~\eqref{rec:omega}, we now have
\begin{align*}
\omega_{n,k-1,k}&=\sum_{p=1}^{\left\lfloor \frac{s}{2}\right\rfloor} \sum_{q=1}^{s+1-2p} \alpha_{s-1}(p,q)\cdot \omega_{n-s,k+s-1-p,k+1-q}
-\sum_{p=1}^{\left\lceil \frac{s}{2}\right\rceil} \sum_{q=1}^{s+2-2p} \alpha_{s}(p,q)\cdot \omega_{n-s+1,k+s-1-p,k+1-q}
\\ &=\sum_{p=1}^{\left\lfloor \frac{s}{2}\right\rfloor} \sum_{q=1}^{s+1-2p} \alpha_{s-1}(p,q)\cdot \omega_{n-s,k+s-1-p,k+1-q}
- \sum_{p=1}^{\left\lceil \frac{s}{2}\right\rceil} \sum_{q=1}^{s+2-2p} \alpha_{s}(p,q)\cdot \omega_{n-s,k+s-p,k+1-q}
\\ &\quad -\sum_{p=1}^{\left\lceil \frac{s}{2}\right\rceil} \sum_{q=1}^{s+2-2p} \alpha_{s}(p,q)\cdot \big( -(s+q-p)\omega_{n-s, k+s-p, k-q}-\omega_{n-s-1, k+s-p, k+1-q} \big)
\\ &= \sum_{p=1}^{\left\lfloor \frac{s+1}{2}\right\rfloor} \sum_{q=1}^{s+2-2p} \alpha_{s}(p,q)\cdot \omega_{n-s-1, k+s-p, k+1-q}
+\sum_{p=1}^{\left\lfloor \frac{s}{2}\right\rfloor} \sum_{q=1}^{s+1-2p} \alpha_{s-1}(p,q)\cdot \omega_{n-s,k+s-1-p,k+1-q}
\\ & \quad -\sum_{p=1}^{\left\lceil \frac{s}{2}\right\rceil} \sum_{q=1}^{s+2-2p} \alpha_{s}(p,q)\cdot \omega_{n-s,k+s-p,k+1-q}
+\sum_{p=1}^{\left\lceil \frac{s}{2}\right\rceil} \sum_{q=1}^{s+2-2p} \alpha_{s}(p,q)\cdot (s+q-p)\omega_{n-s, k+s-p, k-q},
\end{align*}
where we use the fact $\lceil\frac{s}{2}\rceil=\lfloor\frac{s+1}{2}\rfloor$ in the last equality.
Notice that all $\omega$ terms now have the common form $\omega_{n-s,k+s-p,k+1-q}$ after shifting the indices
$p$ and $q$ appropriately. The proof strategy is to show that this identity holds coefficient-wise,
i.e., the coefficients of $\omega_{n-s,k+s-p,k+1-q}$ on both sides are equal. After shifting the indices,
this reduces to the following identity:
\begin{align}
& \sum_{p=2}^{\left\lfloor \frac{s}{2}\right\rfloor+1} \sum_{q=1}^{s+3-2p} \alpha_{s-1}(p-1,q)\cdot \omega_{n-s,k+s-p,k+1-q}
 -\sum_{p=1}^{\left\lceil \frac{s}{2}\right\rceil} \sum_{q=1}^{s+2-2p} \alpha_{s}(p,q)\cdot \omega_{n-s,k+s-p,k+1-q}  \nonumber
\\ &\quad +\sum_{p=1}^{\left\lceil \frac{s}{2}\right\rceil} \sum_{q=2}^{s+3-2p} \alpha_{s}(p,q-1)\cdot (s+q-p-1)\omega_{n-s, k+s-p, k+1-q}\label{Equatio-Need-omega}
\\= &-\sum_{p=1}^{\left\lceil \frac{s+1}{2}\right\rceil} \sum_{q=1}^{s+3-2p} \alpha_{s+1}(p,q)\cdot \omega_{n-s,k+s-p,k+1-q}. \nonumber
\end{align}

We begin by verifying the case $p=1$, that is,  Eq.~\eqref{Equatio-Need-omega} reduces to
\begin{align*}
&-\sum_{q=1}^{s}\alpha_s(1,q) \omega_{n-s,k+s-1,k+1-q}+\sum_{q=2}^{s+1}\alpha_s(1,q-1)(s+q-2)\omega_{n-s,k+s-1,k+1-q}
\\ =&-\sum_{q=1}^{s+1}\alpha_{s+1}(1,q)\omega_{n-s,k+s-1,k+1-q}.
\end{align*}
This is divided into the following three subcases:
\begin{enumerate}
\item When $q=1$, the equality $-\alpha_s(1,1)=-\alpha_{s+1}(1,1)$ follows from~\eqref{Equation-Alpha}.
\item When $q=s+1$, the identity $\alpha_s(1,s)\cdot (2s-1)=-\alpha_{s+1}(1,s+1)$ holds by~\eqref{Equation-Alpha}.
\item When $2\leq q\leq s$, the equality $-\alpha_s(1,q)+\alpha_s(1,q-1)\cdot (s+q-2)=-\alpha_{s+1}(1,q)$ still follows from~\eqref{Equation-Alpha}.
\end{enumerate}

We now turn to the case $p>1$. This requires us to distinguish the parity of $s$. We begin by assuming $s$ is even.

When $p$ attains its maximum value of $\frac{s+2}{2}$, we must have $q=1$, and the identity
\[
\alpha_{s-1}\left(\frac{s}{2},1\right)=-\alpha_{s+1}\left(\frac{s+2}{2},1\right)
\]
follows directly from~\eqref{Equation-Alpha}.

For $2\leq p\leq \frac{s}{2}$, it suffices to prove~\eqref{Equation-Alph-p=2tos}.
\begin{align}
& \sum_{q=1}^{s+3-2p} \alpha_{s-1}(p-1,q)\cdot \omega_{n-s,k+s-p,k+1-q}
 -\sum_{q=1}^{s+2-2p} \alpha_{s}(p,q)\cdot \omega_{n-s,k+s-p,k+1-q} \nonumber
\\ &\quad +\sum_{q=2}^{s+3-2p} \alpha_{s}(p,q-1)\cdot (s+q-p-1)\omega_{n-s, k+s-p, k+1-q} \label{Equation-Alph-p=2tos}
\\= &- \sum_{q=1}^{s+3-2p} \alpha_{s+1}(p,q)\cdot \omega_{n-s,k+s-p,k+1-q}.\nonumber
\end{align}
The verification splits naturally into three cases according to the value of $q$:

\begin{enumerate}
\item If $q=1$, then the identity reduces to
\[
\alpha_{s-1}(p-1,1)-\alpha_s(p,1)=-\alpha_{s+1}(p,1),
\]
which is an immediate consequence of~\eqref{Equation-Alpha}.

\item If $q=s+3-2p$ (the maximum possible value), the required identity becomes
\[
\alpha_{s-1}(p-1,s+3-2p)+\alpha_s(p,s+2-2p)\cdot (2s-3p+2)=-\alpha_{s+1}(p,s+3-2p),
\]
which again follows from~\eqref{Equation-Alpha}.

\item For the interior range $2\leq q\leq s+2-2p$, the identity simplifies to
\[
\alpha_{s-1}(p-1,q)-\alpha_s(p,q)+\alpha_s(p,q-1)\cdot (s+q-p-1)=-\alpha_{s+1}(p,q),
\]
also a direct consequence of~\eqref{Equation-Alpha}.
\end{enumerate}

This completes the proof for even $s$.

Now consider the case where $s$ is odd.
The argument is entirely analogous, with only the boundary values of $p$ adjusted.
When $p=\frac{s+1}{2}$ (the maximum), we have $q\in\{1,2\}$, and the two identities
\begin{align*}
\alpha_{s-1}\left(\frac{s-1}{2},1\right)-\alpha_s\left(\frac{s+1}{2},1\right)&=-\alpha_{s+1}\left(\frac{s+1}{2},1\right),\\
\alpha_{s-1}\left(\frac{s-1}{2},2\right)+\alpha_s\left(\frac{s+1}{2},1\right)\cdot \frac{s+1}{2}
&=-\alpha_{s+1}\left(\frac{s+1}{2},2\right)
\end{align*}
both follow from~\eqref{Equation-Alpha}.
For $2\leq p\leq \frac{s-1}{2}$, the same identity~\eqref{Equation-Alph-p=2tos} holds, and the three-case verification by $q$ proceeds exactly as in the even case.
We omit the repetitive details.

We have thus established that~\eqref{Equatio-Need-omega} holds and that the formula for $\omega_{n,k-1,k}$ is valid for the current value of $s$.
This completes the induction.
\end{proof}

\end{document}
